\documentclass[
  paper=a4,
  tate,
  fontsize=12pt
]{jlreq}

\usepackage{luatexja}
\usepackage{luatexja-fontspec}

% Choose any installed Japanese font.
% HaranoAji fonts are included with TeX Live.
\setmainjfont{HaranoAjiMincho}

\title{郊外の子が市へ}
\author{賈　康麗}
\date{}

\begin{document}

\maketitle

私にとって、課外活動以外に夏休みは新しい経験を得る機会です。八月に私のおばさんがはじめてアメリカに来たので、私たちが町を案内されました。

町の中心部にランニングトラック、バスケットコート、テニスコート、サッカー場、野球場があり、自転車道で自転車が乗れ、図書館で本が読めます。町には小さな酪農場があり、そこで新鮮な肉や乳製品などを買うことができます。酪農場にはアイスクリーム店もあります。ただ地元のアイスクリーム店なのに、複数の賞を受賞することがあります。すべてが小さな町で近くに位置しているので、とても便利です。かなり小さい町ですが、静かで落ち着いて成長できる場所だと思います。

数日後、家族と一緒にニューヨークへ旅行しました。都市では裏庭の木が輝いている高楼に変わってしまいました。マンハッタンの道で、私は熱帯雨林で迷っているアリのように感じました。様々な国籍の人々がどこでもいました。ウォール街やオキュラスなどの賑やかな場所もありました。しかし、他の区は二十世紀に建てられたので、古くて汚いです。地下鉄の駅に入ってすぐに、息を止めなければなりません。蒸し暑い空気の中で、そこのにおいがひどいものでした。都市では、喫煙者もかなり多かったです。ニューヨークには多種多様の人がいるんじゃないか。

そして、通りに路上販売人から弁護士までも、芸術家から大学生までもいました。「ニューヨークには本当に誰にでも空間がある。ここに行ったら、どんな人になりたいだろうか。」と考えました。市に行き、すぐに社会の底辺に落ちたら、大変でしょう。それを防ぐために、私は今から頑張るべきです。

今日、私は「あ、夏休みがもう終わったかな。」と思わずに、新たなやる気でいっぱいで、勉強する準備ができました！　将来のために、とりあえず、今は知識を広げる時間です。

\end{document}